% Generated draft: requires mathematical and source review.
\documentclass[11pt]{article}
\usepackage[margin=25mm]{geometry}
\usepackage{fontspec}
\setmainfont{Latin Modern Roman}
\usepackage{amsmath,amssymb,mathtools,unicode-math}
\setmathfont{Latin Modern Math}
\usepackage{enumitem,needspace,xurl,hyperref}
\setlength{\parindent}{0pt}
\setlength{\parskip}{5pt}
\setlength{\emergencystretch}{3em}
\newcommand{\N}{\mathbb{N}}
\newcommand{\Z}{\mathbb{Z}}
\newcommand{\Q}{\mathbb{Q}}
\newcommand{\R}{\mathbb{R}}
\newcommand{\C}{\mathbb{C}}
\newcommand{\F}{\mathbb{F}}
\newcommand{\A}{\mathbb{A}}
\newcommand{\PP}{\mathbb{P}}
\newcommand{\E}{\mathbb{E}}
\newcommand{\eps}{\varepsilon}
\newcommand{\dd}{\,\mathrm{d}}
\begin{document}
\section*{UnsolvedMath ID: 30000287}
\textbf{Source problem code:} catalogue entry
\subsection{Definitions and mathematical statement}
Let \(\omega:\Z\to(0,\infty)\) satisfy, for some finite \(c\geq1\),
\[
c^{-1}\leq\frac{\omega(n+1)}{\omega(n)}\leq c\qquad(n\in\Z).
\]
Use the weighted Hilbert-space convention
\[
\ell^2_\omega(\Z)=\left\{a:\Z\to\C:\|a\|_\omega^2=\sum_{n\in\Z}|a(n)|^2\omega(n)^2<\infty\right\}.
\]
Define the bilateral translation \((Ta)(n)=a(n-1)\). The ratio bounds make \(T\) bounded and invertible, with \(\|T\|,\|T^{-1}\|\leq c\); translation invariance of the ambient space does not mean that translations are isometries.

On the usual all-integer-translations reading, the question is whether for every such \(\omega\) there exists a closed complex linear subspace \(Y\) such that
\[
\{0\}\subsetneq Y\subsetneq\ell^2_\omega(\Z),\qquad
T^kY\subseteq Y\quad\text{for every }k\in\Z.
\]
Equivalently, \(TY=Y\). Closedness is essential to the usual invariant-subspace question. Invariance only under \(T\), with no condition on \(T^{-1}\), is a weaker and distinct formulation; the report's precise use of translation-invariant should be confirmed.

Some conventions put \(\omega(n)\) rather than \(\omega(n)^2\) in the norm. The definition above makes the chosen convention explicit; replacing the weight by its square root translates between the two, preserving the class of positive bounded-ratio weights, though not the name of each individual weight. No symmetry, subexponential growth, or extra regularity is imposed.

The supplied assessment checked 2026-08-21 reports that the universal existence question remains open, with positive results for important subclasses of admissible weights. The subclasses and exact hypotheses are not reproduced.
\subsection{Short English statement}
Does every positive bounded-ratio weight on \(\Z\) give a weighted square-summable sequence space with a nonzero proper closed subspace invariant under all integer translations? The norm and two-sided invariance conventions are stated explicitly.
\subsection{Sources}
\begin{itemize}
\item Local UnsolvedMath registry, permanent ID 30000287, catalogue entry, \emph{Invariant Subspaces of Weighted Translation Operators}.
\item \emph{Mini-Workshop: Operators on Spaces of Analytic Functions} (2006), Oberwolfach Reports, pp. 2043--2090, DOI \url{https://doi.org/10.4171/owr/2005/36}.
\item Supplied reference labeled \emph{Translation invariant subspaces of weighted Hilbert spaces}: \url{https://congreso.us.es/ceacyto/2006/pdfs/esterle.pdf}.
\item Supplied literature triage checked 2026-08-21 and source-verification claim of comparison with the official report.
\end{itemize}
\end{document}
